\chapter{Selected Source Code Listings}
\label{app:source-code-listings}

This appendix provides critical source code listings that implement the core mathematical and algorithmic mechanisms discussed throughout the manuscript. The listings are derived from the Wolverine ecosystem (\texttt{wolverine-sih}, \texttt{wolverine-db}) and the visual interaction layer (\texttt{Vzeya}). 

\section{Graph Projection via Bounded Breadth-First Search}
\label{sec:code-projector}

As discussed in the architectural chapters, the Wolverine system explicitly avoids database-level recursive mechanisms, such as PostgreSQL recursive Common Table Expressions (CTEs), for graph projection. Instead, it relies on a highly optimized, strictly in-memory Bounded Breadth-First Search (BFS) implemented in \texttt{projector.ts}. This approach allows for tighter control over traversal depths and memory consumption during real-time network analytics.

\begin{lstlisting}[language=TypeScript, caption={In-memory Bounded BFS for graph projection (\texttt{projector.ts})}, label={lst:projector}]
export class GraphProjector {
  /**
   * Executes an in-memory Bounded BFS to project the local neighborhood
   * of a given entity without relying on database recursive CTEs.
   */
  public projectNeighborhood(
    startNodeId: string,
    maxDepth: number = 3
  ): Graph {
    const visited = new Set<string>();
    const queue: Array<{ id: string; depth: number }> = [];
    const resultGraph = new Graph();

    queue.push({ id: startNodeId, depth: 0 });
    visited.add(startNodeId);

    while (queue.length > 0) {
      const { id, depth } = queue.shift()!;
      const node = this.nodeCache.get(id);

      if (node) {
        resultGraph.addNode(node);
      }

      if (depth < maxDepth) {
        const edges = this.edgeCache.getEdges(id);
        for (const edge of edges) {
          resultGraph.addEdge(edge);
          const neighborId = edge.target === id ? edge.source : edge.target;

          if (!visited.has(neighborId)) {
            visited.add(neighborId);
            queue.push({ id: neighborId, depth: depth + 1 });
          }
        }
      }
    }

    return resultGraph;
  }
}
\end{lstlisting}

\section{Entity Resolution and Linear Temporal Decay}
\label{sec:code-entity-resolver}

The \texttt{entity\_resolver.ts} module demonstrates the implementation of the probabilistic matching algorithms utilized within the Wolverine ecosystem. Notably, it incorporates a linear 30-day temporal decay function to adjust confidence scores, explicitly eschewing exponential decay models to guarantee mathematical predictability over the resolution lifecycle. Furthermore, the system enforces strict deterministic thresholds: a strong match threshold of 0.92 and a candidate threshold of 0.70.

\begin{lstlisting}[language=TypeScript, caption={Entity resolution logic with linear temporal decay and strict thresholds (\texttt{entity\_resolver.ts})}, label={lst:entity-resolver}]
export class EntityResolver {
  private static readonly STRONG_MATCH_THRESHOLD = 0.92;
  private static readonly CANDIDATE_THRESHOLD = 0.70;
  private static readonly MAX_DECAY_DAYS = 30;

  /**
   * Calculates the final confidence score using a linear 30-day temporal decay.
   * Exponential decay is explicitly avoided to maintain predictable 
   * auditability in resolution scores.
   */
  private applyLinearTemporalDecay(
    baseScore: number, 
    lastSeenDate: Date, 
    currentDate: Date
  ): number {
    const diffTime = Math.abs(currentDate.getTime() - lastSeenDate.getTime());
    const diffDays = Math.ceil(diffTime / (1000 * 60 * 60 * 24));
    
    if (diffDays >= EntityResolver.MAX_DECAY_DAYS) {
      return 0.0; // Complete decay after 30 days
    }
    
    // Linear decay penalty
    const decayFactor = 1.0 - (diffDays / EntityResolver.MAX_DECAY_DAYS);
    return baseScore * decayFactor;
  }

  public resolveEntities(entityA: Entity, entityB: Entity): ResolutionResult {
    const baseScore = this.calculateSimilarity(entityA, entityB);
    const finalScore = this.applyLinearTemporalDecay(
      baseScore, 
      entityB.lastSeen, 
      new Date()
    );

    if (finalScore >= EntityResolver.STRONG_MATCH_THRESHOLD) {
      return { status: 'MERGE', confidence: finalScore };
    } else if (finalScore >= EntityResolver.CANDIDATE_THRESHOLD) {
      return { status: 'REVIEW', confidence: finalScore };
    }
    
    return { status: 'REJECT', confidence: finalScore };
  }
}
\end{lstlisting}

\section{CSS-Based CRT Post-Processing}
\label{sec:code-crt}

Within the Vzeya visual interaction layer, the Cathode-Ray Tube (CRT) terminal effect is achieved purely through Cascading Style Sheets (CSS). The implementation actively avoids computationally expensive WebGL GLSL shaders. Instead, it utilizes the \texttt{repeating-linear-gradient} property to synthesize scanlines, ensuring lightweight cross-browser compatibility and significantly lower computational overhead.

\begin{lstlisting}[language=JavaScript, caption={CSS-based scanline generation for CRT effects (\texttt{CRTPostProcessing.tsx})}, label={lst:crt-css}]
import React from 'react';
import './CRTStyles.css';

/**
 * Renders a CRT overlay using CSS repeating-linear-gradient.
 * WebGL/GLSL shaders are completely avoided in Vzeya for this effect.
 */
export const CRTPostProcessing: React.FC<{ children: React.ReactNode }> = ({ 
  children 
}) => {
  return (
    <div className="crt-container">
      <div className="crt-scanlines" style={{
        background: `
          linear-gradient(
            rgba(18, 16, 16, 0) 50%, 
            rgba(0, 0, 0, 0.25) 50%
          ), 
          repeating-linear-gradient(
            90deg, 
            rgba(255, 0, 0, 0.06), 
            rgba(0, 255, 0, 0.02), 
            rgba(0, 0, 255, 0.06) 3px
          )`
        ,
        backgroundSize: '100% 2px, 3px 100%',
        pointerEvents: 'none'
      }} />
      <div className="crt-content">
        {children}
      </div>
    </div>
  );
};
\end{lstlisting}
